\clearpage
\pdfbookmark[1]{Reports. 30 July 1849.}{anzeige_1849_122_124}
\SourceLeaf{122}
\par\noindent\hfill\rule{0.30\linewidth}{0.4pt}\hfill\mbox{}\par
\begin{center}
\textbf{Reports of the Royal Society of Sciences at Göttingen. 30 July 1849.}
\end{center}

\par\noindent\hfill\rule{0.30\linewidth}{0.4pt}\hfill\mbox{}\par

The lecture given by Privy Court Councillor \textsc{Gauss} at the meeting of the Royal Society of Sciences on 16 July concerned the
\begin{center}
\textit{Theory of algebraic equations.}
\end{center}
We give here only the following report of the detailed memoir, which will shortly appear in print.

It consists of two divisions, each devoted to a separate subject.

The first division deals with the fundamental theorem in the theory of algebraic equations, namely the factorization of every polynomial in one variable, on which the author had already published a memoir in 1799 under the title \textit{A new proof of the theorem that every polynomial} \textit{in one variable can be resolved into real factors of the first or second degree.} This memoir had a twofold purpose: first, to show that all the proofs of the stated theorem attempted up to that time were inadequate and illusory; and second, to give a fully rigorous new proof. It would be unnecessary to return to the first subject. It must, however, be noted that the author later followed that first new proof with two others, which are found in the 3rd volume
\SourceLeaf{123}
of our Society's \textit{Commentationes recentiores}, and that a fourth was first put forward by \textsc{Cauchy}. All four of these fully rigorous proofs rest on as many different foundations, but agree in that they all \textit{directly} establish only the existence of one factor of the function in question. This does not detract from the rigour of the proofs. For it is clear that, after this one factor has been removed from the given function, a similar function of lower degree remains, to which the theorem can be applied again, and that continued repetition of the procedure will finally yield a complete factorization of the original function into factors of the stated kind. Nevertheless, every proof evidently attains greater completeness if it is shown to be capable of making the existence of all the individual factors directly evident at once. The first proof does indeed have this property, as was already indicated, though not developed further, in the memoir mentioned above (Art. 23). The present paper now gives this extension, but as part of a reworking of the memoir of 1799 rather than as an addition to it. The proof is presented here afresh, independently, but with some essential changes in its course that yield a considerable simplification. Its whole clothing has also been changed. In the earlier work this had deliberately been chosen so that any involvement of imaginary quantities could be avoided entirely. The argument, however, is most closely connected with the conception of complex quantities, both in its \textit{origin} and in its true nature. Attentive readers of that memoir have not failed to notice this since the conception became established in mathematics. The author has therefore considered it appropriate now to abandon that first clothing and return the argument to its proper field, whose basic concepts are now familiar to everyone, especially since this yields not only greater simplicity but also somewhat greater generality.

The second division is devoted to algebraic equations with three terms. These have the distinctive feature that some of the methods for the numerical solution of equations can achieve a flexibility and elegance in their case far beyond what their application to equations of less simple form allows. This applies in particular to solution
\SourceLeaf{124}
by infinite series and to the indirect method. The development of these methods for equations of that form therefore seems to deserve a separate treatment all the more because such equations in fact occur very frequently.

The author has, however, entirely excluded solution by infinite series from his present purpose, remarking only that for \textit{every} root of such an equation, whether real or imaginary, a convergent series proceeding according to an easily recognized law can be found. Beautiful as this manner of solution is from a general theoretical standpoint, in actual practical applications the indirect methods will nevertheless be preferred in all cases where that convergence is not very rapid.

These indirect methods are the subject of the second division. There are \textit{two} methods here, since the procedure required to determine the imaginary roots differs entirely from that to be used for the real roots. On other occasions the author has already given a few examples of the latter in particular cases, while also indicating the general applicability of the procedure. Although this generalization presents no difficulty at all, the author nevertheless believed that readers would welcome a collection of instructions ready for use, especially since the number of cases to be distinguished and treated individually is considerable.

Finding the imaginary roots by an indirect method is much more difficult than finding the real roots, unless somewhat approximate values are already known from elsewhere: the former must be sought in an infinite region of two dimensions, the latter only in an infinite region of one dimension. For equations with three terms, this difficulty can be avoided by a simple and very natural device that, however, seems not to have been used before. For this problem too, the instructions required for its solution are given in full and ready for use.
\par\noindent\hfill\rule{0.30\linewidth}{0.4pt}\hfill\mbox{}\par
